\documentclass[12pt,a4paper]{article}
\usepackage[utf8]{inputenc}
\usepackage[T1]{fontenc}
\usepackage[french]{babel}
\usepackage{amsmath,amssymb}
\usepackage{graphicx}
\usepackage{float}
\usepackage{geometry}
\geometry{margin=2.5cm}

\title{Analyse en Composantes Principales}
\author{Votre Nom}
\date{\today}

\begin{document}

\maketitle

\begin{abstract}
Résumé de votre analyse en composantes principales.
\end{abstract}

\section{Introduction}

L'analyse en composantes principales (ACP) est une méthode statistique...

\section{Méthodologie}

\subsection{Principes mathématiques}

L'ACP cherche à réduire la dimensionnalité des données...

\subsection{Algorithme}

Les étapes principales de l'ACP sont :
\begin{enumerate}
    \item Centrer et réduire les données
    \item Calculer la matrice de covariance
    \item Déterminer les valeurs propres et vecteurs propres
    \item Sélectionner les composantes principales
\end{enumerate}

\section{Application numérique}

\subsection{Jeu de données}

Présentez vos données...

\subsection{Résultats}

Analysez les résultats obtenus...

\section{Conclusion}

Synthétisez vos findings...

\begin{thebibliography}{9}
\bibitem{exemple}
Auteur. \textit{Titre}. Éditeur, année.
\end{thebibliography}

\end{document}